%! Designing Experiments — Unit 3, Lesson 3.4 — Student Packet Cover Sheet
\documentclass[12pt]{article}
\usepackage{apstats-article}
\usepackage{apstats-boxes}
\usepackage{ltablex}
\keepXColumns

\begin{document}

% Full-bleed navy banner. \coverbanner MEASURES the title block and sizes the
% band to it, so a lesson title long enough to wrap grows the band rather than
% running out the bottom of it.
\coverbanner{Unit 3}{Lesson 3.4 \quad Designing Experiments}

% NAMESTRIP: this is the ONLY component that carries the name row.
\namedateperiod
\vspace{0.18in}

\boxguard[14]
\begin{learningtargetbox}
\small
% GRADUAL RELEASE: the vocabulary is taught in the guided notes, so the targets name it
% plainly. The AP Learning Objectives (VAR-3.A--D) stay in the teacher-facing plan.
\begin{enumerate}[leftmargin=1.5em, itemsep=5pt]
  \item \ldots{} name the \textbf{experimental units}, \textbf{factors}, \textbf{treatments},
        and \textbf{response variable} of an experiment.
  \item \ldots{} explain the four principles of a well-designed experiment ---
        \textbf{comparison}, \textbf{random assignment}, \textbf{replication}, and
        \textbf{control} --- including a \textbf{control group}, a \textbf{placebo}, and
        \textbf{blinding}.
  \item \ldots{} choose between a \textbf{completely randomized design}, a
        \textbf{randomized block design}, and a \textbf{matched pairs design}, and explain
        why chance must still assign the treatments inside every block or pair.
\end{enumerate}
\end{learningtargetbox}

\vspace{0.14in}

\boxguard[14]
\begin{tocbox}
\small
% A CONTENTS LIST ONLY: no Score column and no Total row. Whether a component counts is said in
% its description: AP Practice is extra credit, scored 1 to 5; the Homework is scored.
\renewcommand{\arraystretch}{1.5}
\begin{tabularx}{\linewidth}{@{} c l X @{}}
\rowcolor{sky}
\textbf{\#} & \textbf{Component} & \textbf{Description} \\
\midrule
1 & Warm-Up & Breakfast and GPAs, a review game in two class periods, and splitting a class by
                 chance \\
2 & Guided Notes \& Practice & Gum and memory, a sports drink blocked by runner, a shoe sole
                 tested on both feet, and a zinc lozenge for colds \\
3 & AP Practice & Tomato plants, two sunscreens, and a math-practice app: AP-style multiple
                 choice and one free-response set --- extra credit, scored 1 to 5 \\
4 & Homework & Corn on two kinds of field, a new pillow, and a joint chew for older dogs ---
                 \textbf{scored 1 to 5, due the first class after two study halls} \\
\end{tabularx}
\end{tocbox}

\vspace{0.14in}

\boxguard[8]
\begin{remindbox}
\small
An experiment \textbf{does something} to its units and watches what happens, so it can show
cause and effect --- but only if \textbf{chance} decides who gets which treatment. Blocking
and pairing put alike units side by side so the comparison is sharper; they do \emph{not}
replace the coin flip. \textbf{Block on something that affects the response}, then randomize
inside every block. A placebo and blinding keep what people expect from posing as a treatment
effect.
\end{remindbox}

\end{document}
